\documentclass[12pt, a4paper]{article}
\usepackage[utf8]{inputenc}
\usepackage{amsmath, amsthm, amssymb, amsfonts}
\usepackage{geometry}
\geometry{top=2.5cm, bottom=2.5cm, left=2.5cm, right=2.5cm}
\usepackage{hyperref}
\usepackage{cite}

% Document structure
\newtheorem{proposition}{Proposition}[section]
\newtheorem{remark}{Remark}[section]

\title{\textbf{An Exploratory Survey on Reformulations, Operator Models, and the Off-Axis Cross-Pairing Barrier for the Riemann Hypothesis}}
\author{\textbf{Antigravity Research Collaborative} \\ \small In Pair-Programming Partnership with the User}
\date{August 2026}

\begin{document}

\maketitle

\begin{abstract}
This working paper provides an exploratory survey of several modern equivalent reformulations of the Riemann Hypothesis (RH) across operator theory, harmonic analysis, and complex function theory. We review the Weil quadratic criterion, continuous screw-kernel representations (Suzuki 2026), and half-plane logarithmic derivative behaviors. Through detailed structural analysis, we highlight the central analytical bottleneck common to all these frameworks: establishing the non-negativity of off-axis bilinear pairings $-\sum \widehat{v}(\rho)\overline{\widehat{v}(1-\overline{\rho})}$ without pre-supposing the Riemann Hypothesis. We also document why standard complex-analytic continuation tools from $\operatorname{Re}(s) > 1$ cannot bypass interior pole exclusions, and contrast this structure with the explicit metric breakdown observed in Epstein zeta functions.
\end{abstract}

\tableofcontents

\section{Introduction and Scope}
The Riemann Hypothesis (RH), asserting that all non-trivial zeros of $\zeta(s)$ lie on $\operatorname{Re}(s) = 1/2$, possesses numerous equivalent reformulations in modern mathematics. Notable frameworks include:
\begin{itemize}
    \item \textbf{Harmonic Analysis}: André Weil's positivity criterion (1952) \cite{Weil1952}.
    \item \textbf{Noncommutative Geometry}: Alain Connes' adelic trace formula (1999) \cite{Connes1999} and spectral triples (2025) \cite{CCM2025}.
    \item \textbf{Canonical Hamiltonian Systems}: M. Suzuki's continuous screw-kernel operators (2026) \cite{Suzuki2026}.
    \item \textbf{Geometric Function Theory}: Half-plane Carath\'eodory / Herglotz logarithmic derivative positivity \cite{deBranges1968}.
\end{itemize}

The purpose of this document is not to claim a proof of RH, but rather to map out these parallel mathematical languages, verify their conceptual connections, and precisely isolate the shared mathematical barrier where all current techniques stall.

\section{The Weil Quadratic Functional}
Let $\mathcal{T}_{\text{Weil}}$ be the admissible space of smooth test functions $g(x) \in C_c^\infty(\mathbb{R})$ satisfying decay and pole-neutralization conditions $\widehat{g}(0) = 0, \widehat{g}(\pm i/2) = 0$. The Weil quadratic functional is given by:
\begin{equation}
Q_W(v) = W(v * \widetilde{v}), \quad \widetilde{v}(x) = \overline{v(-x)},
\end{equation}
with the functional decomposition:
\begin{equation}
W(f) = W_{\text{arch}}(f) + W_{\text{pole}}(f) - \sum_{n=1}^\infty \frac{\Lambda(n)}{\sqrt{n}} \left( f(\log n) + f(-\log n) \right).
\end{equation}

\begin{remark}[The Arithmetic Sign]
For any positive-type convolution $g_v = v * \widetilde{v}$, the arithmetic matrix on discrete Galerkin truncations satisfies $M_{\text{prime}} \succeq 0$. However, because this term is subtracted in the explicit formula ($-M_{\text{prime}} \preceq 0$), prime contributions energetically subtract from the positive Archimedean background, meaning $M_{\text{prime}} \succeq 0$ alone does not imply $Q_W \ge 0$.
\end{remark}

\section{Survey of Five Equivalent Perspectives}
Several modern frameworks describe the same underlying condition in different mathematical vernaculars:

\begin{enumerate}
    \item \textbf{Operator Theoretic}: The continuum limit operator $D_\infty = \lim_{a\to\infty} A_a$ constructed from the localized screw kernel is self-adjoint with strictly real spectrum \cite{Suzuki2026}.
    \item \textbf{Complex Analytic}: The logarithmic derivative $\operatorname{Re}(\xi'/\xi(s)) > 0$ strictly in the open half-plane $\mathbb{H}_{1/2} = \{\operatorname{Re}(s) > 1/2\}$.
    \item \textbf{Harmonic Analytic}: $Q_W(g) \ge 0$ for all admissible test functions $g \in \mathcal{T}_{\text{Weil}}$ \cite{Weil1952}.
    \item \textbf{Analytic Number Theoretic}: The prime sum $\sum \frac{\Lambda(n)}{n^\sigma}\cos(t\log n)$ is uniformly dominated by the Archimedean background $\frac{1}{2}\log(|t|/2\pi) + \mathcal{O}(1)$ across $\sigma > 1/2$.
    \item \textbf{Geometric Function Theoretic}: The Schwarz--Pick ratio $\mathcal{S}(\sigma, t) = \frac{(\sigma-1/2)|F'|}{\operatorname{Re}(F)} \le 1$ holds throughout $\mathbb{H}_{1/2}$.
\end{enumerate}

These formulations rephrase RH in different functional settings, but none provide an independent bypass of the arithmetic zero-distribution.

\section{The Core Bottleneck: Off-Axis Cross-Pairing}
Decomposing the explicit formula over the zero set $Z_\zeta$, the quadratic zero contribution is:
\begin{equation}
W_{\text{zeros}}(v * \widetilde{v}) = -\sum_{\rho \in Z_\zeta} \widehat{v}(\rho) \overline{\widehat{v}(1 - \overline{\rho})}.
\end{equation}

\begin{itemize}
    \item \textbf{Under RH ($\beta = 1/2$)}: Since $1 - \overline{\rho} = \rho$, the pairing collapses to $-|\widehat{v}(\rho)|^2 \le 0$, resulting in a sum of non-negative modulus squares.
    \item \textbf{Without assuming RH ($\beta = 1/2 + \delta, \delta \ne 0$)}: The cross-pairing $\widehat{v}(\rho) \overline{\widehat{v}(1 - \overline{\rho})}$ involves differing exponential factors $e^{\pm \delta x}$ and cannot be simplified to a modulus square.
\end{itemize}

Proving that $Q_W(v) \ge 0$ unconditionally without assuming $\delta = 0$ is mathematically identical to solving the Riemann Hypothesis itself.

\section{Analytical Limitations of Geometric Continuation}
In the absolute convergence half-plane $\operatorname{Re}(s) > 1$, the Dirichlet series $\sum \frac{\Lambda(n)}{n^s}$ provides explicit bounds. However:
\begin{itemize}
    \item Schwarz--Pick contraction $\mathcal{S}(s) \le 1$ on a subdomain $\operatorname{Re}(s) > 1$ cannot be analytically continued into $\operatorname{Re}(s) > 1/2$ to rule out interior poles.
    \item Meromorphic continuation does not preserve positive-real-part cones in the presence of unexcluded zeros.
\end{itemize}

\section{Epstein Zeta Counterexamples}
For Epstein zeta functions $Z_E(s)$ associated with binary quadratic forms with class number $h > 1$, off-critical zeros $\rho_E = \beta_E + i\gamma_E$ ($\beta_E > 1/2$) exist. In this setting:
\begin{enumerate}
    \item The pseudo-prime coefficients $a_q$ take negative values, breaking scalar Gram decompositions.
    \item For $\sigma < \beta_E$, the field $\operatorname{Re}(Z_E'/Z_E)$ plunges below zero, entering a negative phase.
    \item The Schwarz--Pick ratio $\mathcal{S}_E \to +\infty$ as $\sigma \to \beta_E$, demonstrating metric breakdown.
\end{enumerate}
This demonstrates that purely geometric frameworks must inherently depend on the multiplicative uniqueness of prime numbers (the Euler product) to prevent metric degeneration.

\section{Conclusion}
This survey highlights the structural unity among modern operator, harmonic, and geometric formulations of the Riemann Hypothesis. All paths converge onto the same fundamental obstacle: controlling the sign of off-axis bilinear pairings against the arithmetic explicit formula. Recognizing this shared barrier clarifies the precise limitations of current methods and provides a disciplined foundation for future research.

\begin{thebibliography}{99}
\bibitem{Weil1952}
A. Weil, \emph{Sur les ``formules explicites'' de la th\'eorie des nombres premiers}, Comm. S\'em. Math. Univ. Lund [Medd. Lunds Univ. Mat. Sem.] (1952), 252--265.

\bibitem{Connes1999}
A. Connes, \emph{Trace formula in noncommutative geometry and the zeros of the Riemann zeta function}, Selecta Math. (N.S.) \textbf{5} (1999), no. 1, 29--106.

\bibitem{deBranges1968}
L. de Branges, \emph{Hilbert Spaces of Entire Functions}, Prentice-Hall, Englewood Cliffs, N.J., 1968.

\bibitem{CCM2025}
A. Connes, C. Consani, H. Moscovici, \emph{The Zeta Spectral Triple and Logarithmic Derivatives of the Riemann Zeta Function}, arXiv:2511.22755 (2025).

\bibitem{Suzuki2026}
M. Suzuki, \emph{Weil's quadratic form via the screw function}, arXiv:2606.09096 (2026).

\bibitem{Groskin2026a}
A. Groskin, \emph{High-Precision Approximation of Riemann Zeros via the Truncated Weil Form}, arXiv:2605.20224 (2026).
\end{thebibliography}

\end{document}
